\section{Introduction}

Prognostics and Health Management (PHM) for high-value rotating machinery, such as commercial aviation turbofan engines, is critical to modern industrial operations~\cite{saxena2008damage,huang2021survey}. Operational maintenance strategies have traditionally navigated a severe trade-off between corrective maintenance and time-based preventive overhauls. Run-to-failure corrective strategies maximize the physical utilization of machine life but risk uncontained in-service failures, introducing severe flight delays and catastrophic safety hazards. Conversely, scheduled overhauls at fixed flight intervals prevent unexpected shutdowns but scrap substantial components with significant Remaining Useful Life (RUL) intact~\cite{shehadeh2024evaluating}.

Predictive Maintenance (PdM) leverages multi-sensor telemetry to estimate RUL dynamically. Although modern deep neural networks and gradient boosting regressors achieve low point-prediction errors, point estimates $\hat{y}$ provide no measure of predictive uncertainty. In sequential replacement decisions, treating deterministic point estimates as ground truth is operationally hazardous: overestimating remaining life leads to in-service breakdown, while underestimating remaining life precipitates premature replacement~\cite{chen2022data,diao2026turbofan}.

To enable mathematically reliable decision boundaries, prognostic models must produce calibrated prediction intervals $C(x) = [L(x), U(x)]$ targeting nominal coverage $1-\alpha$. Conformal prediction~\cite{javanmardi2023conformal} provides a distribution-free framework for uncertainty quantification.

IntelliTwin uses an engine-weighted conformal construction motivated by hierarchical conformal prediction. Formal coverage results require the hierarchical sampling and exchangeability assumptions stated in the underlying theory. Equal engine weighting does not establish these assumptions for ordered, lifetime-stratified degradation trajectories. This study therefore evaluates marginal and subgroup coverage empirically. It does not claim guaranteed coverage for each engine, lifecycle stage, operating regime, complete trajectory, or adaptively selected maintenance time. Nominal 90% coverage is a target, not a guaranteed realized coverage proportion. Maintenance outcomes are evaluated separately under the frozen simulator. Furthermore, downstream replacement decisions depend primarily on the lower prediction bound $L(x)$, creating an asymmetric operational penalty that symmetric calibration methods fail to address~\cite{cortesgomez2025utility,robinson2026riskaware}.

The IntelliTwin project investigates \textit{decision-aware, condition-adaptive conformal calibration} on turbofan degradation benchmarks, seeking to align conditional prediction-interval reliability with sequential maintenance utility. The system is evaluated as an offline, trace-driven prognostic decision-support research prototype; physical engine integration and operational airline validation have not been demonstrated, and maintenance costs are expressed in Simulated Cost Units (SCU) representing synthetic benchmark preferences rather than validated airline economics.
